\epigraph{It is virtually impossible to allow marginal productivity theory to determine the optimal level of capital, and then have marginal efficiency of investment theory determine the optimal level of investment without thereby eliminating the flow investment term entirely.}{\textit{Trygve Haavelmo} (1960)}

This chapter introduces adjustment costs through the CES transformation technology. The key insight is that the elasticity of transformation between consumption and investment plays a role analogous to the elasticity of substitution in consumer theory. When this elasticity is finite, the economy faces implicit adjustment costs: the marginal rate of transformation depends on the allocation, making the effective interest rate endogenous.

We develop the theory progressively: from a two-period model that introduces Tobin's $Q$ and compares with standard adjustment cost formulations, through finite-horizon dynamics, to economies with multiple capital types and incomplete depreciation. Throughout, the CES structure delivers clean results: zero profits in decentralization, wealth effects operating only through prices, and tractable aggregation.

% ==============================================================
\section{The Two-Period Problem with Adjustment Costs}
\label{sec:ac-two-period}
% ==============================================================

We now introduce adjustment costs through the CES aggregator. The key insight is that the same machinery governing substitution across goods also governs substitution between consumption and investment. When the aggregator has finite elasticity, transforming output into consumption and investment is costly---the economy faces an implicit adjustment cost even without explicit convex costs in the investment technology.

% --------------------------------------------------------------
\subsection{From Linear to CES Transformation}
\label{subsec:ac-linear-to-ces}
% --------------------------------------------------------------

In the standard neoclassical model, output $Y$ splits linearly into consumption $C$ and investment $I$:
\begin{equation}
C + I = Y.
\label{eq:ac-linear-constraint}
\end{equation}
This formulation assumes perfect substitutability: one unit of output can become one unit of consumption or one unit of investment at the margin, with no loss. The relative ``price'' of investment in terms of consumption is always unity.

\paragraph{Introducing the CES Transformation.}
Now suppose that converting output into consumption and investment requires an aggregation technology. Output $Y$ is related to consumption and investment through a CES aggregator:
\begin{equation}
\Hfun(C, I; \bom, \psi) = Y,
\label{eq:ac-ces-constraint}
\end{equation}
where $\Hfun$ is a CES function with elasticity parameter $\psi$ and weights $\bom = (\omega_C, \omega_I)$ satisfying $\omega_C + \omega_I = 1$. In canonical form:
\begin{equation}
\Hfun(C, I) = \left( \omega_C^{1-\psi} C^{\psi} + \omega_I^{1-\psi} I^{\psi} \right)^{1/\psi}.
\label{eq:ac-H-canonical}
\end{equation}

The parameter $\psi$ governs the elasticity of transformation between consumption and investment. Define the elasticity of substitution:
\begin{equation}
\eta \equiv \frac{1}{1 - \psi}.
\label{eq:ac-eta-def}
\end{equation}

\begin{center}
\begin{tabular}{@{}lccc@{}}
\toprule
Case & Power $\psi$ & Elasticity $\eta$ & Interpretation \\
\midrule
Linear & $\psi = 1$ & $\eta = \infty$ & Perfect substitutes \\
Cobb-Douglas & $\psi = 0$ & $\eta = 1$ & Unit elasticity \\
Leontief & $\psi \to -\infty$ & $\eta = 0$ & Perfect complements \\
Intermediate & $\psi \in (0,1)$ & $\eta > 1$ & High substitutability \\
Intermediate & $\psi < 0$ & $\eta < 1$ & Low substitutability \\
\bottomrule
\end{tabular}
\end{center}

When $\eta < \infty$, the economy faces an implicit adjustment cost: shifting resources from consumption to investment (or vice versa) is costly because the transformation frontier is curved.

\paragraph{The Production Possibility Frontier.}
Given output $Y$, the constraint $\Hfun(C, I) = Y$ defines a \emph{production possibility frontier} (PPF) in $(C, I)$ space. Solving for $I$ as a function of $C$:
\begin{equation}
I = \left( \frac{Y^{\psi} - \omega_C^{1-\psi} C^{\psi}}{\omega_I^{1-\psi}} \right)^{1/\psi}.
\label{eq:ac-ppf}
\end{equation}

The slope of the PPF is the \emph{marginal rate of transformation} (MRT):
\begin{equation}
\text{MRT} \equiv -\frac{dI}{dC} = \frac{\omega_C^{1-\psi}}{\omega_I^{1-\psi}} \left( \frac{C}{I} \right)^{\psi - 1} = \frac{\omega_C^{1-\psi}}{\omega_I^{1-\psi}} \left( \frac{I}{C} \right)^{1-\psi}.
\label{eq:ac-mrt}
\end{equation}

When $\psi < 1$ (i.e., $\eta < \infty$), the MRT depends on the allocation $(C, I)$. The cost of producing an additional unit of investment---in terms of foregone consumption---varies along the frontier. This is the essence of adjustment costs.

\begin{calloutnote}[The Price of Investment]
The MRT can be interpreted as the relative price of investment in terms of consumption. When $\eta < \infty$, this price is endogenous: it depends on how much the economy is investing. High investment rates make additional investment expensive at the margin.

This connects to the classic $q$-theory of investment: the shadow price of capital (Tobin's $q$) reflects the marginal cost of transforming output into installed capital.
\end{calloutnote}

% --------------------------------------------------------------
\subsection{The Two-Period Setup}
\label{subsec:ac-two-period-setup}
% --------------------------------------------------------------

Consider a two-period economy, $t \in \{0, 1\}$. The household enters period 0 with capital stock $k_0$ and chooses consumption and investment to maximize lifetime utility.

\paragraph{Technology.}
Production is linear in capital with productivity $A$:
\begin{equation}
Y_t = A k_t.
\label{eq:ac-production}
\end{equation}
Capital fully depreciates ($\delta = 1$), so next-period capital equals current investment:
\begin{equation}
k_1 = I_0.
\label{eq:ac-capital-accum}
\end{equation}
The CES aggregator governs the transformation of output:
\begin{equation}
\Hfun(C_t, I_t) = Y_t = A k_t.
\label{eq:ac-resource}
\end{equation}

\paragraph{Preferences.}
The household has time-separable preferences with discount factor $\beta \in (0,1)$ and period utility $u(C) = \frac{C^{1-1/\theta} - 1}{1 - 1/\theta}$, where $\theta > 0$ is the intertemporal elasticity of substitution (IES):
\begin{equation}
U = u(C_0) + \beta \, u(C_1).
\label{eq:ac-utility}
\end{equation}

\paragraph{Terminal Condition.}
In the final period, there is no investment ($I_1 = 0$), so all output is consumed:
\begin{equation}
C_1 = Y_1 = A k_1 = A I_0.
\label{eq:ac-terminal}
\end{equation}
Note that in period 1, the CES constraint becomes $\Hfun(C_1, 0) = Y_1$. For the canonical CES with $\psi < 1$, this requires $C_1 = (\omega_C^{1-\psi})^{-1/\psi} Y_1$ when $I_1 = 0$. To avoid this complication, we assume the CES constraint applies only to period 0; in period 1, output converts directly to consumption.

% --------------------------------------------------------------
\subsection{The Optimization Problem}
\label{subsec:ac-optimization}
% --------------------------------------------------------------

The household's problem is:
\begin{equation}
\max_{C_0, I_0} \left\{ u(C_0) + \beta \, u(A I_0) \right\} \quad \text{subject to} \quad \Hfun(C_0, I_0) = A k_0.
\label{eq:ac-problem}
\end{equation}

This is a CES optimization problem, but with a twist: the constraint is nonlinear (a CES aggregator rather than a linear budget). We can apply the transformation approach from earlier chapters.

\paragraph{The Dual Interpretation.}
View the problem as choosing the allocation $(C_0, I_0)$ on the PPF to maximize a CES objective in $(C_0, C_1)$ where $C_1 = A I_0$. The objective is:
\begin{equation}
\max_{C_0, I_0} \; \V(C_0, A I_0; \tilde{\bom}, \rho) \quad \text{subject to} \quad \Hfun(C_0, I_0; \bom, \psi) = Y_0,
\label{eq:ac-dual-problem}
\end{equation}
where $\V$ is the intertemporal aggregator with power $\rho = 1 - 1/\theta$ and weights $(\tilde{\omega}_0, \tilde{\omega}_1)$ derived from $\beta$ as in Section~\ref{sec:two-period}.

This is a \emph{CES-on-CES} problem: a CES objective subject to a CES constraint. The interaction between the two elasticities---$\theta$ (intertemporal) and $\eta$ (transformation)---determines the equilibrium.

% --------------------------------------------------------------
\subsection{First-Order Conditions and the Euler Equation}
\label{subsec:ac-foc}
% --------------------------------------------------------------

\paragraph{The Lagrangian.}
Form the Lagrangian:
\begin{equation}
\mathcal{L} = u(C_0) + \beta \, u(A I_0) + \lambda \left[ Y_0 - \Hfun(C_0, I_0) \right].
\label{eq:ac-lagrangian}
\end{equation}

\paragraph{First-Order Conditions.}
Differentiating with respect to $C_0$ and $I_0$:
\begin{align}
\frac{\partial \mathcal{L}}{\partial C_0} &= u'(C_0) - \lambda \frac{\partial \Hfun}{\partial C_0} = 0, \label{eq:ac-foc-c} \\
\frac{\partial \mathcal{L}}{\partial I_0} &= \beta A \, u'(A I_0) - \lambda \frac{\partial \Hfun}{\partial I_0} = 0. \label{eq:ac-foc-i}
\end{align}

The partial derivatives of the CES aggregator are:
\begin{align}
\frac{\partial \Hfun}{\partial C_0} &= \omega_C^{1-\psi} C_0^{\psi-1} \Hfun^{1-\psi} = \omega_C^{1-\psi} \left( \frac{C_0}{\Hfun} \right)^{\psi-1}, \label{eq:ac-dH-dC} \\
\frac{\partial \Hfun}{\partial I_0} &= \omega_I^{1-\psi} I_0^{\psi-1} \Hfun^{1-\psi} = \omega_I^{1-\psi} \left( \frac{I_0}{\Hfun} \right)^{\psi-1}. \label{eq:ac-dH-dI}
\end{align}

\paragraph{The Modified Euler Equation.}
Taking the ratio of \cref{eq:ac-foc-c,eq:ac-foc-i}:
\begin{equation}
\frac{u'(C_0)}{\beta A \, u'(C_1)} = \frac{\partial \Hfun / \partial C_0}{\partial \Hfun / \partial I_0} = \frac{\omega_C^{1-\psi}}{\omega_I^{1-\psi}} \left( \frac{C_0}{I_0} \right)^{\psi-1}.
\label{eq:ac-euler-ratio}
\end{equation}

The right-hand side is the MRT from \cref{eq:ac-mrt}. Define the \emph{implicit price of investment}:
\begin{equation}
q \equiv \text{MRT} = \frac{\omega_C^{1-\psi}}{\omega_I^{1-\psi}} \left( \frac{C_0}{I_0} \right)^{\psi-1}.
\label{eq:ac-q-def}
\end{equation}

The Euler equation becomes:
\begin{equation}
\boxed{u'(C_0) = \beta \, q \, A \, u'(C_1).}
\label{eq:ac-euler}
\end{equation}

\begin{calloutimportant}[The Endogenous Interest Rate]
Compare with the standard Euler equation $u'(C_0) = \beta R \, u'(C_1)$. The gross return $R$ is replaced by $q \cdot A$, where:
\begin{itemize}[nosep]
\item $A$ is the marginal product of capital (the ``technological'' return).
\item $q$ is the relative price of investment (the MRT).
\end{itemize}

The effective interest rate is $R = q \cdot A$. When $\eta = \infty$ (linear transformation), $q = 1$ and $R = A$---the standard case. When $\eta < \infty$, $q$ depends on the investment rate, making the interest rate endogenous.
\end{calloutimportant}

% --------------------------------------------------------------
\subsection{Characterizing the Solution}
\label{subsec:ac-solution}
% --------------------------------------------------------------

\paragraph{Solving for the Allocation.}
With CRRA utility, $u'(C) = C^{-1/\theta}$, the Euler equation \cref{eq:ac-euler} becomes:
\begin{equation}
C_0^{-1/\theta} = \beta \, q \, A \, C_1^{-1/\theta} = \beta \, q \, A \, (A I_0)^{-1/\theta}.
\label{eq:ac-euler-crra}
\end{equation}

Rearranging:
\begin{equation}
\frac{C_1}{C_0} = (\beta \, q \, A)^{\theta}.
\label{eq:ac-consumption-ratio}
\end{equation}

Using $C_1 = A I_0$ and $q = \frac{\omega_C^{1-\psi}}{\omega_I^{1-\psi}} (C_0/I_0)^{\psi-1}$:
\begin{equation}
\frac{A I_0}{C_0} = \left( \beta A \cdot \frac{\omega_C^{1-\psi}}{\omega_I^{1-\psi}} \left( \frac{C_0}{I_0} \right)^{\psi-1} \right)^{\theta}.
\label{eq:ac-ratio-expanded}
\end{equation}

Let $x \equiv I_0 / C_0$ denote the investment-to-consumption ratio. Then:
\begin{equation}
A x = \left( \beta A \cdot \frac{\omega_C^{1-\psi}}{\omega_I^{1-\psi}} \cdot x^{1-\psi} \right)^{\theta}.
\label{eq:ac-x-equation}
\end{equation}

Taking logs:
\begin{equation}
\log(A x) = \theta \left[ \log(\beta A) + (1-\psi) \log\left( \frac{\omega_C}{\omega_I} \right) + (1-\psi) \log x \right].
\label{eq:ac-log-equation}
\end{equation}

Collecting terms in $\log x$:
\begin{equation}
\log x \left[ 1 - \theta(1-\psi) \right] = \theta \log(\beta A) + \theta(1-\psi) \log\left( \frac{\omega_C}{\omega_I} \right) - \log A.
\label{eq:ac-log-x-solve}
\end{equation}

Define:
\begin{equation}
\xi \equiv 1 - \theta(1-\psi) = 1 - \frac{\theta}{\eta}.
\label{eq:ac-xi-def}
\end{equation}

The parameter $\xi$ captures the interaction between the IES $\theta$ and the transformation elasticity $\eta$. When $\xi \neq 0$:
\begin{equation}
\log x = \frac{1}{\xi} \left[ (\theta - 1) \log A + \theta \log \beta + \theta(1-\psi) \log\left( \frac{\omega_C}{\omega_I} \right) \right].
\label{eq:ac-log-x-solution}
\end{equation}

\paragraph{Special Cases.}

\medskip
\noindent
\textit{Case 1: Linear transformation ($\eta = \infty$, $\psi = 1$).} Then $\xi = 1$ and $q = \omega_C / \omega_I$ (constant). The solution reduces to the standard consumption-savings problem with exogenous interest rate $R = A \omega_C / \omega_I$.

\medskip
\noindent
\textit{Case 2: Cobb-Douglas transformation ($\eta = 1$, $\psi = 0$).} Then $\xi = 1 - \theta$. The solution is:
\begin{equation}
\log x = \frac{(\theta - 1) \log A + \theta \log \beta + \theta \log(\omega_C / \omega_I)}{1 - \theta}.
\label{eq:ac-cobb-douglas}
\end{equation}
For $\theta = 1$, this requires separate analysis (the denominator vanishes).

\medskip
\noindent
\textit{Case 3: Unit IES ($\theta = 1$).} Then $\xi = \psi$ and:
\begin{equation}
\log x = \frac{\log \beta + (1-\psi) \log(\omega_C / \omega_I)}{\psi}.
\label{eq:ac-unit-ies}
\end{equation}
The investment rate is independent of $A$---income and substitution effects cancel.

\begin{calloutnote}[The Role of $\xi$]
The parameter $\xi = 1 - \theta/\eta$ determines how the investment rate responds to changes in productivity $A$:
\begin{itemize}[nosep]
\item When $\xi > 0$ (i.e., $\theta < \eta$): Higher $A$ increases investment if $\theta > 1$ and decreases investment if $\theta < 1$.
\item When $\xi < 0$ (i.e., $\theta > \eta$): The responses are reversed.
\item When $\xi = 0$ (i.e., $\theta = \eta$): The investment rate is independent of $A$.
\end{itemize}
The interaction between intertemporal substitution and transformation elasticity creates rich comparative statics.
\end{calloutnote}

% --------------------------------------------------------------
\subsection{The Transformation Approach}
\label{subsec:ac-transformation}
% --------------------------------------------------------------

We now solve the same problem using the transformation approach, which reveals the structural role of the adjustment cost.

\paragraph{Change of Variables.}
Define the \emph{output shares}:
\begin{equation}
s_C \equiv \frac{\omega_C^{1-\psi} C_0^{\psi}}{Y_0^{\psi}}, \qquad s_I \equiv \frac{\omega_I^{1-\psi} I_0^{\psi}}{Y_0^{\psi}}.
\label{eq:ac-shares-def}
\end{equation}

By the CES constraint, $s_C + s_I = 1$. These shares measure the ``contribution'' of each component to total output, analogous to expenditure shares in the consumer problem.

\paragraph{Inverting the Transformation.}
From \cref{eq:ac-shares-def}:
\begin{equation}
C_0 = \left( \frac{s_C}{\omega_C^{1-\psi}} \right)^{1/\psi} Y_0, \qquad I_0 = \left( \frac{s_I}{\omega_I^{1-\psi}} \right)^{1/\psi} Y_0.
\label{eq:ac-invert}
\end{equation}

\paragraph{Transformed Objective.}
Substituting into the utility function and using $C_1 = A I_0$:
\begin{align}
U &= u\left( \left( \frac{s_C}{\omega_C^{1-\psi}} \right)^{1/\psi} Y_0 \right) + \beta \, u\left( A \left( \frac{s_I}{\omega_I^{1-\psi}} \right)^{1/\psi} Y_0 \right) \notag \\
&= u\left( \left( \frac{s_C}{\omega_C^{1-\psi}} \right)^{1/\psi} Y_0 \right) + \beta \, u\left( \left( \frac{s_I}{\omega_I^{1-\psi}} \right)^{1/\psi} A Y_0 \right).
\label{eq:ac-transformed-U}
\end{align}

This is a standard two-good problem in the shares $(s_C, s_I)$ with the linear constraint $s_C + s_I = 1$. The ``prices'' are embedded in the transformation from shares to quantities.

\paragraph{Effective Prices.}
Define the effective price of consumption in period $t$ (in terms of the output share):
\begin{equation}
\tilde{p}_C = \left( \omega_C^{1-\psi} \right)^{1/\psi}, \qquad \tilde{p}_I = \left( \omega_I^{1-\psi} \right)^{1/\psi} / A.
\label{eq:ac-effective-prices}
\end{equation}

The ratio $\tilde{p}_I / \tilde{p}_C$ captures the relative cost of future consumption (via investment) in terms of current consumption.

% --------------------------------------------------------------
\subsection{Tobin's $Q$ and the Price of Investment}
\label{subsec:ac-tobin-q}
% --------------------------------------------------------------

The marginal rate of transformation $q$ defined in \cref{eq:ac-q-def} has a natural interpretation as \emph{Tobin's $Q$}---the shadow price of installed capital in terms of consumption goods.

\begin{definition}[Tobin's $Q$]
\label{def:tobin-q}
Tobin's $Q$ is the marginal cost of investment in terms of consumption:
\begin{equation}
Q \equiv -\frac{dC}{dI} \bigg|_{\Hfun = Y} = \frac{\partial \Hfun / \partial I}{\partial \Hfun / \partial C} = \frac{\omega_I^{1-\psi}}{\omega_C^{1-\psi}} \left( \frac{C}{I} \right)^{1-\psi}.
\label{eq:ac-tobin-q}
\end{equation}
Equivalently, $Q = 1/q$ where $q$ is the MRT defined earlier.
\end{definition}

Note the convention: $q = \text{MRT} = -dI/dC$ measures consumption foregone per unit of investment, while $Q = 1/q = -dC/dI$ measures the cost of investment in consumption units. In the $q$-theory literature, it is common to use $Q$ (or lowercase $q$) for the latter. We follow this convention henceforth.

\paragraph{Interpretation.}
When $Q > 1$, investment is ``expensive''---one unit of installed capital costs more than one unit of output. When $Q < 1$, investment is ``cheap.'' The deviation of $Q$ from unity reflects the curvature of the transformation frontier.

In the CES formulation:
\begin{equation}
Q = \frac{\omega_I^{1-\psi}}{\omega_C^{1-\psi}} \left( \frac{C}{I} \right)^{1-\psi} = \frac{\omega_I^{1-\psi}}{\omega_C^{1-\psi}} \left( \frac{C}{I} \right)^{1/\eta}.
\label{eq:ac-Q-ces}
\end{equation}

When $\eta < \infty$, $Q$ depends on the consumption-investment ratio. High investment rates (low $C/I$) imply high $Q$: the marginal unit of investment becomes increasingly costly.

\begin{calloutnote}[The Euler Equation with Tobin's $Q$]
The Euler equation \cref{eq:ac-euler} can be rewritten as:
\begin{equation}
u'(C_0) = \beta \frac{A}{Q} u'(C_1).
\label{eq:ac-euler-Q}
\end{equation}
The effective gross return is $R = A/Q$: the marginal product of capital $A$ divided by the cost of acquiring that capital $Q$. This is the fundamental insight of $q$-theory: investment is governed by the ratio of marginal benefit ($A$) to marginal cost ($Q$).
\end{calloutnote}

% --------------------------------------------------------------
\subsection{Comparison with Standard Adjustment Costs}
\label{subsec:ac-comparison}
% --------------------------------------------------------------

The literature on investment typically specifies adjustment costs in a different form. We now compare our CES formulation with the standard approach, highlighting several advantages of the CES specification.

\paragraph{The Standard Formulation.}
The canonical adjustment cost specification in macroeconomics (following Lucas and Prescott, 1971; Hayashi, 1982) is:
\begin{equation}
C + I \left[ 1 + \frac{\phi}{2} \left( \frac{I}{K} \right) \right] = Y,
\label{eq:ac-standard}
\end{equation}
or equivalently, defining the adjustment cost function $\Phi(I, K) = \frac{\phi}{2} (I/K)^2 K$:
\begin{equation}
C + I + \Phi(I, K) = Y, \qquad \Phi(I, K) = \frac{\phi}{2} \frac{I^2}{K}.
\label{eq:ac-standard-phi}
\end{equation}

This specification has become standard because it delivers Tobin's $Q$ as a sufficient statistic for investment and because it is tractable in continuous time.

\paragraph{The Odd Normalization.}
The standard formulation raises an immediate question: \emph{why normalize by capital $K$?} The adjustment cost $\Phi(I, K) = \frac{\phi}{2} I^2 / K$ is homogeneous of degree one in $(I, K)$, which ensures that the value function is linear in $K$ along a balanced growth path. But this normalization is not derived from first principles---it is imposed for tractability.

Several issues arise:
\begin{enumerate}[label=(\roman*), nosep]
\item \textbf{Arbitrary scale}: Why should adjustment costs scale with existing capital? If adjustment costs reflect installation, learning, or organizational frictions, there is no obvious reason why doubling the capital stock should double the cost of a given investment level.

\item \textbf{Asymmetry}: The cost depends on $I/K$, making it asymmetric between investment and consumption. There is no corresponding friction on the consumption side.

\item \textbf{Singularity at $K = 0$}: When $K \to 0$, the adjustment cost per unit of investment explodes: $\Phi/I = \frac{\phi}{2} I/K \to \infty$. This creates technical difficulties for economies starting from low capital.
\end{enumerate}

\paragraph{The CES Alternative.}
The CES formulation $\Hfun(C, I) = Y$ avoids these issues:
\begin{enumerate}[label=(\roman*), nosep]
\item \textbf{Symmetric treatment}: Both consumption and investment enter the transformation technology symmetrically. The friction applies to the \emph{composition} of output, not to investment per se.

\item \textbf{No arbitrary normalization}: The curvature is governed by a single parameter $\eta$ (the elasticity of transformation) with clear economic meaning.

\item \textbf{Well-behaved limits}: As $\eta \to \infty$, we recover the frictionless case $C + I = Y$. As $\eta \to 0$, the technology becomes Leontief.

\item \textbf{Independence from $K$}: The transformation frontier depends on output $Y$, not on the capital stock. This is natural if the friction reflects the technology for converting final goods into consumption and investment goods.
\end{enumerate}

\begin{calloutcaution}[When Does $K$ Belong in Adjustment Costs?]
The dependence on $K$ in the standard formulation can be motivated if adjustment costs represent \emph{internal} installation costs---frictions that occur within the firm and scale with firm size. In contrast, the CES formulation is natural when adjustment costs represent \emph{external} transformation costs---frictions in converting output across uses, perhaps reflecting sectoral reallocation or the technology of the capital goods industry.

The appropriate specification depends on the economic context. For aggregate analysis where ``adjustment costs'' proxy for various frictions, the CES formulation offers a parsimonious and well-behaved alternative.
\end{calloutcaution}

% --------------------------------------------------------------
\subsection{Decentralization and the Zero-Profit Property}
\label{subsec:ac-decentralization}
% --------------------------------------------------------------

A key advantage of the CES formulation is that it decentralizes cleanly with \emph{zero profits}. The standard adjustment cost formulation, by contrast, generates profits that must be allocated to shareholders.

\paragraph{Decentralizing the CES Economy.}
Consider a competitive firm that purchases output $Y$ at price $1$ (the numeraire) and sells consumption goods at price $p_C$ and investment goods at price $p_I$. The firm's problem is:
\begin{equation}
\max_{C, I} \; p_C C + p_I I \quad \text{subject to} \quad \Hfun(C, I) \leq Y.
\label{eq:ac-firm-problem}
\end{equation}

The first-order conditions require:
\begin{equation}
\frac{p_C}{p_I} = \frac{\partial \Hfun / \partial C}{\partial \Hfun / \partial I} = \frac{\omega_C^{1-\psi}}{\omega_I^{1-\psi}} \left( \frac{C}{I} \right)^{\psi - 1}.
\label{eq:ac-firm-foc}
\end{equation}

\paragraph{Zero Profits.}
The CES aggregator is homogeneous of degree one: $\Hfun(\lambda C, \lambda I) = \lambda \Hfun(C, I)$. By Euler's theorem:
\begin{equation}
\Hfun(C, I) = \frac{\partial \Hfun}{\partial C} C + \frac{\partial \Hfun}{\partial I} I.
\label{eq:ac-euler-theorem}
\end{equation}

At the optimum, marginal products equal prices (up to the Lagrange multiplier $\mu$):
\begin{equation}
p_C = \mu \frac{\partial \Hfun}{\partial C}, \qquad p_I = \mu \frac{\partial \Hfun}{\partial I}.
\label{eq:ac-marginal-pricing}
\end{equation}

Substituting into revenue:
\begin{equation}
p_C C + p_I I = \mu \left( \frac{\partial \Hfun}{\partial C} C + \frac{\partial \Hfun}{\partial I} I \right) = \mu \, \Hfun(C, I) = \mu Y.
\label{eq:ac-revenue}
\end{equation}

For the constraint to bind, $\mu = 1$ (the shadow price of output equals its market price). Therefore:
\begin{equation}
\boxed{\text{Revenue} = p_C C + p_I I = Y = \text{Cost}.}
\label{eq:ac-zero-profit}
\end{equation}

\begin{calloutimportant}[Zero Profits from Constant Returns]
The zero-profit result follows from constant returns to scale in the transformation technology. This is the same logic as in standard production theory: with CRS and competitive pricing, firms earn zero profits.

This property is attractive for general equilibrium analysis. There are no profits to distribute, no need to specify ownership shares, and the decentralization is clean. Households simply face prices $(p_C, p_I)$ for the two goods.
\end{calloutimportant}

\paragraph{Contrast with Standard Adjustment Costs.}
In the standard formulation \cref{eq:ac-standard-phi}, the firm's problem is:
\begin{equation}
\max_{I} \; p_I I - I - \Phi(I, K) = \max_{I} \; (p_I - 1) I - \frac{\phi}{2} \frac{I^2}{K}.
\label{eq:ac-standard-firm}
\end{equation}

The first-order condition gives:
\begin{equation}
p_I = 1 + \phi \frac{I}{K} \equiv Q.
\label{eq:ac-standard-Q}
\end{equation}

Here $Q = p_I$ is Tobin's $Q$---the market price of installed capital. The firm's profit is:
\begin{equation}
\Pi = (p_I - 1) I - \frac{\phi}{2} \frac{I^2}{K} = \phi \frac{I}{K} \cdot I - \frac{\phi}{2} \frac{I^2}{K} = \frac{\phi}{2} \frac{I^2}{K} > 0.
\label{eq:ac-standard-profit}
\end{equation}

\emph{Profits are strictly positive whenever $I > 0$.} These profits must be distributed to shareholders, requiring a specification of ownership. This complicates the decentralization and raises questions about the incidence of adjustment costs.

\paragraph{The Spurious Wealth Effect.}
The positive profits in the standard formulation create a conceptual problem beyond mere accounting. In the CES economy, the household's budget constraint is:
\begin{equation}
C + Q \cdot I = Y.
\label{eq:ac-ces-budget}
\end{equation}
All wealth effects operate through prices: changes in $Q$ affect the household's opportunity set, but total resources available equal output $Y$. This is the standard benchmark in consumer theory---wealth effects are fully summarized by the budget constraint evaluated at market prices.

In the standard adjustment cost economy, the household's budget constraint becomes:
\begin{equation}
C + Q \cdot I = Y + \Pi(I, K), \qquad \Pi = \frac{\phi}{2} \frac{I^2}{K}.
\label{eq:ac-standard-budget}
\end{equation}
Now there is an \emph{additional} wealth effect: profits $\Pi$ depend on the investment rate $I/K$. When investment rises, profits rise, expanding the household's budget. This creates a feedback loop that is absent in the frictionless model and absent in the CES formulation.

\begin{calloutimportant}[Wealth Effects and the Price System]
In the CES formulation, all information relevant to the household's decision is summarized by prices $(p_C, p_I)$ and output $Y$. The household can read off wealth effects directly from the budget constraint---no additional channel operates.

In the standard formulation, wealth effects have two components:
\begin{enumerate}[nosep]
\item The direct effect through prices (as in the CES case).
\item An indirect effect through profits, which depend on equilibrium quantities.
\end{enumerate}

This indirect channel breaks the clean separation between prices and quantities that characterizes competitive equilibrium. The household's wealth depends not just on prices but on the \emph{level} of investment---a quantity that the household itself influences.

For welfare analysis, this distinction matters: the incidence of adjustment costs differs across the two formulations. In the CES economy, costs are fully reflected in relative prices. In the standard economy, part of the cost is offset by profit income, with the net burden depending on ownership shares.
\end{calloutimportant}

\paragraph{Market Prices in the CES Economy.}
From the firm's first-order conditions and the zero-profit condition, we can solve for equilibrium prices. Normalizing $p_C = 1$ (consumption as numeraire):
\begin{equation}
p_I = \frac{\partial \Hfun / \partial I}{\partial \Hfun / \partial C} = \frac{\omega_I^{1-\psi}}{\omega_C^{1-\psi}} \left( \frac{C}{I} \right)^{1-\psi} = Q.
\label{eq:ac-p-I}
\end{equation}

The price of investment goods equals Tobin's $Q$. When $C/I$ is high (low investment rate), $p_I = Q$ is low---investment goods are cheap. When $C/I$ is low (high investment rate), $p_I = Q$ is high---investment goods are expensive.

The household's budget constraint is:
\begin{equation}
C + Q \cdot I = Y,
\label{eq:ac-hh-budget}
\end{equation}
which is exactly the CES constraint $\Hfun(C, I) = Y$ expressed in value terms at equilibrium prices.

\begin{table}[H]
\centering
\caption{Comparison of Adjustment Cost Formulations}
\label{tab:ac-comparison}
\renewcommand{\arraystretch}{1.5}
\begin{tabular}{@{}lcc@{}}
\toprule
\textbf{Property} & \textbf{Standard: $\Phi = \frac{\phi}{2} I^2/K$} & \textbf{CES: $\Hfun(C,I) = Y$} \\
\midrule
Normalization & By capital $K$ (ad hoc) & None required \\[4pt]
Symmetry & Asymmetric (only investment) & Symmetric \\[4pt]
Profits & $\Pi = \frac{\phi}{2} I^2/K > 0$ & $\Pi = 0$ \\[4pt]
Wealth effects & Through prices \emph{and} profits & Through prices only \\[4pt]
Tobin's $Q$ & $Q = 1 + \phi I/K$ & $Q = \dfrac{\omega_I^{1-\psi}}{\omega_C^{1-\psi}} \left( \dfrac{C}{I} \right)^{1-\psi}$ \\[10pt]
Behavior at $K = 0$ & Singular & Well-defined \\[4pt]
Curvature parameter & $\phi$ & $\eta = 1/(1-\psi)$ \\[4pt]
\bottomrule
\end{tabular}
\end{table}

% --------------------------------------------------------------
\subsection{Summary}
\label{subsec:ac-two-period-summary}
% --------------------------------------------------------------

\begin{table}[H]
\centering
\caption{The Two-Period Problem with CES Adjustment Costs}
\label{tab:ac-two-period-summary}
\renewcommand{\arraystretch}{1.6}
\begin{tabular}{@{}ll@{}}
\toprule
\textbf{Object} & \textbf{Formula} \\
\midrule
CES constraint & $\Hfun(C, I) = \left( \omega_C^{1-\psi} C^{\psi} + \omega_I^{1-\psi} I^{\psi} \right)^{1/\psi} = Y$ \\[6pt]
Transformation elasticity & $\eta = 1/(1-\psi)$ \\[6pt]
Tobin's $Q$ & $Q = \dfrac{\omega_I^{1-\psi}}{\omega_C^{1-\psi}} \left( \dfrac{C}{I} \right)^{1-\psi}$ \\[10pt]
Euler equation & $u'(C_0) = \beta \dfrac{A}{Q} u'(C_1)$ \\[8pt]
Effective interest rate & $R = A/Q$ (endogenous) \\[6pt]
Key parameter & $\xi = 1 - \theta/\eta$ \\[6pt]
Zero-profit condition & $C + Q \cdot I = Y$ \\[6pt]
\bottomrule
\end{tabular}
\end{table}

\begin{callouttip}[Key Insights]
The two-period problem with CES adjustment costs reveals that:
\begin{itemize}[nosep]
\item \textbf{Tobin's $Q$} emerges naturally as the marginal rate of transformation---the shadow price of investment in terms of consumption.
\item \textbf{The interest rate is endogenous}: $R = A/Q$ depends on the investment rate through Tobin's $Q$.
\item \textbf{Zero profits}: The CES formulation decentralizes with zero profits, unlike the standard $\phi(I/K)^2$ specification.
\item \textbf{No odd normalization}: The curvature is governed by a single parameter $\eta$ without arbitrary dependence on $K$.
\item \textbf{The interaction parameter} $\xi = 1 - \theta/\eta$ governs comparative statics---the interplay between intertemporal substitution and transformation elasticity.
\end{itemize}
This structure extends naturally to multiple periods and to settings with uncertainty.
\end{callouttip}


% ==============================================================
\section{The Finite-Horizon Problem}
\label{sec:ac-T-period}
% ==============================================================

We now extend the two-period analysis to a finite horizon of $T+1$ periods, $t \in \{0, 1, \ldots, T\}$. This extension reveals the full dynamic structure of adjustment costs: how investment at any date responds to productivity changes at any other date, and how Tobin's $Q$ evolves along the transition path.

% --------------------------------------------------------------
\subsection{Setup}
\label{subsec:ac-T-setup}
% --------------------------------------------------------------

\paragraph{Technology.}
Production is linear in capital with productivity $A_t$:
\begin{equation}
Y_t = A_t k_t.
\label{eq:ac-T-production}
\end{equation}
Capital fully depreciates ($\delta = 1$), so next-period capital equals current investment:
\begin{equation}
k_{t+1} = I_t, \qquad t = 0, 1, \ldots, T-1.
\label{eq:ac-T-capital}
\end{equation}
The CES aggregator governs the transformation of output in each period:
\begin{equation}
\Hfun(C_t, I_t) = Y_t, \qquad t = 0, 1, \ldots, T-1.
\label{eq:ac-T-resource}
\end{equation}
In the terminal period $T$, there is no investment: $C_T = Y_T$.

\paragraph{Preferences.}
The household has time-separable preferences with discount factor $\beta \in (0,1)$ and period utility $u(C) = \frac{C^{1-1/\theta} - 1}{1 - 1/\theta}$:
\begin{equation}
U = \sum_{t=0}^{T} \beta^t \, u(C_t).
\label{eq:ac-T-utility}
\end{equation}

\paragraph{State and Control Variables.}
The state variable is capital $k_t$. The control variables are consumption $C_t$ and investment $I_t$, linked by the CES constraint. Given the constraint, we can take $I_t$ as the single control: $C_t$ is then determined by $\Hfun(C_t, I_t) = A_t k_t$.

% --------------------------------------------------------------
\subsection{The Recursive Formulation}
\label{subsec:ac-T-recursive}
% --------------------------------------------------------------

\paragraph{The Bellman Equation.}
Define $V_t(k)$ as the value of entering period $t$ with capital $k$. The Bellman equation is:
\begin{equation}
V_t(k) = \max_{I} \left\{ u(C) + \beta V_{t+1}(I) \right\} \quad \text{subject to} \quad \Hfun(C, I) = A_t k,
\label{eq:ac-T-bellman}
\end{equation}
with terminal condition $V_T(k) = u(A_T k)$.

\paragraph{Eliminating $C$.}
From the CES constraint, consumption is implicitly defined by $\Hfun(C, I) = Y$. For the canonical CES:
\begin{equation}
C = \left( \frac{Y^\psi - \omega_I^{1-\psi} I^\psi}{\omega_C^{1-\psi}} \right)^{1/\psi} \equiv C(I; Y).
\label{eq:ac-T-C-implicit}
\end{equation}

The Bellman equation becomes:
\begin{equation}
V_t(k) = \max_{I} \left\{ u\bigl(C(I; A_t k)\bigr) + \beta V_{t+1}(I) \right\}.
\label{eq:ac-T-bellman-reduced}
\end{equation}

\paragraph{First-Order Condition.}
Differentiating with respect to $I$:
\begin{equation}
u'(C) \frac{\partial C}{\partial I} + \beta V_{t+1}'(I) = 0.
\label{eq:ac-T-foc}
\end{equation}

From the implicit function theorem applied to $\Hfun(C, I) = Y$:
\begin{equation}
\frac{\partial C}{\partial I} = -\frac{\partial \Hfun / \partial I}{\partial \Hfun / \partial C} = -\frac{1}{Q},
\label{eq:ac-T-dCdI}
\end{equation}
where $Q$ is Tobin's $Q$ as defined in \cref{eq:ac-tobin-q}. The first-order condition becomes:
\begin{equation}
\frac{u'(C_t)}{Q_t} = \beta V_{t+1}'(k_{t+1}).
\label{eq:ac-T-foc-Q}
\end{equation}

\paragraph{Envelope Condition.}
Differentiating the value function with respect to $k$:
\begin{equation}
V_t'(k) = u'(C) \frac{\partial C}{\partial Y} \cdot A_t = u'(C_t) \cdot \frac{A_t}{Q_t} \cdot \frac{1}{1 - s_I},
\label{eq:ac-T-envelope-raw}
\end{equation}
where $s_I = \omega_I^{1-\psi} I^\psi / Y^\psi$ is the investment share in the CES aggregator.

For the CES function, $\partial C / \partial Y = (C/Y) / (1 - s_I)$ by homogeneity arguments. At the optimum, the envelope condition simplifies. Using the homogeneity of $\Hfun$:
\begin{equation}
V_t'(k) = \frac{u'(C_t)}{Q_t} \cdot A_t.
\label{eq:ac-T-envelope}
\end{equation}

% --------------------------------------------------------------
\subsection{The Euler Equation and Tobin's $Q$}
\label{subsec:ac-T-euler}
% --------------------------------------------------------------

Combining the first-order condition \cref{eq:ac-T-foc-Q} with the envelope condition \cref{eq:ac-T-envelope} (shifted forward one period):
\begin{equation}
\frac{u'(C_t)}{Q_t} = \beta \cdot \frac{u'(C_{t+1})}{Q_{t+1}} \cdot A_{t+1}.
\label{eq:ac-T-euler-combined}
\end{equation}

Rearranging:
\begin{equation}
\boxed{u'(C_t) = \beta \frac{A_{t+1}}{Q_{t+1}} Q_t \, u'(C_{t+1}).}
\label{eq:ac-T-euler}
\end{equation}

\paragraph{Interpretation.}
The Euler equation has a natural interpretation in terms of returns:
\begin{itemize}[nosep]
\item $Q_t$ is the cost of acquiring one unit of capital at date $t$ (in consumption units).
\item $A_{t+1}$ is the marginal product of that capital at date $t+1$.
\item $Q_{t+1}$ is the cost of capital at date $t+1$---relevant if the household could instead acquire capital next period.
\end{itemize}

Define the \emph{return on capital}:
\begin{equation}
R_{t+1} \equiv \frac{A_{t+1}}{Q_{t+1}} \cdot Q_t.
\label{eq:ac-T-return}
\end{equation}

The Euler equation becomes the standard form:
\begin{equation}
u'(C_t) = \beta R_{t+1} \, u'(C_{t+1}).
\label{eq:ac-T-euler-standard}
\end{equation}

\begin{calloutimportant}[The Return with Adjustment Costs]
The return $R_{t+1} = \frac{A_{t+1}}{Q_{t+1}} Q_t$ has two components:
\begin{enumerate}[nosep]
\item \textbf{Dividend yield}: $A_{t+1}/Q_{t+1}$ is the marginal product per unit of capital value.
\item \textbf{Capital gain}: $Q_t/Q_{t+1}$ captures changes in Tobin's $Q$. If $Q$ is falling ($Q_{t+1} < Q_t$), the return is higher---capital acquired cheaply today can be ``sold'' at a higher implicit price tomorrow.
\end{enumerate}

In the frictionless case ($\eta = \infty$), $Q_t = Q_{t+1} = 1$ and $R_{t+1} = A_{t+1}$. With adjustment costs, returns include capital gains/losses on $Q$.
\end{calloutimportant}

% --------------------------------------------------------------
\subsection{The Sequence Formulation}
\label{subsec:ac-T-sequence}
% --------------------------------------------------------------

An alternative to the recursive approach is to solve the problem in sequence space. This formulation connects directly to the Jacobian methods discussed in earlier chapters.

\paragraph{The Planner's Problem.}
The planner maximizes:
\begin{equation}
\max_{\{C_t, I_t, k_{t+1}\}} \sum_{t=0}^{T} \beta^t u(C_t)
\label{eq:ac-T-planner-obj}
\end{equation}
subject to:
\begin{align}
\Hfun(C_t, I_t) &= A_t k_t, \qquad t = 0, \ldots, T-1, \label{eq:ac-T-planner-ces} \\
k_{t+1} &= I_t, \qquad t = 0, \ldots, T-1, \label{eq:ac-T-planner-accum} \\
C_T &= A_T k_T, \label{eq:ac-T-planner-terminal}
\end{align}
with $k_0$ given.

\paragraph{Lagrangian.}
Form the Lagrangian with multipliers $\lambda_t$ on the CES constraints and $\mu_t$ on the capital accumulation equations:
\begin{equation}
\mathcal{L} = \sum_{t=0}^{T} \beta^t u(C_t) + \sum_{t=0}^{T-1} \lambda_t \left[ A_t k_t - \Hfun(C_t, I_t) \right] + \sum_{t=0}^{T-1} \mu_t \left[ I_t - k_{t+1} \right].
\label{eq:ac-T-lagrangian}
\end{equation}

\paragraph{First-Order Conditions.}
\begin{align}
C_t: & \quad \beta^t u'(C_t) = \lambda_t \frac{\partial \Hfun}{\partial C_t}, \qquad t = 0, \ldots, T-1, \label{eq:ac-T-foc-C} \\
I_t: & \quad \lambda_t \frac{\partial \Hfun}{\partial I_t} = \mu_t, \qquad t = 0, \ldots, T-1, \label{eq:ac-T-foc-I} \\
k_{t+1}: & \quad \mu_t = \lambda_{t+1} A_{t+1}, \qquad t = 0, \ldots, T-2, \label{eq:ac-T-foc-k} \\
C_T: & \quad \beta^T u'(C_T) = \mu_{T-1} A_T. \label{eq:ac-T-foc-CT}
\end{align}

\paragraph{Recovering Tobin's $Q$.}
From \cref{eq:ac-T-foc-C,eq:ac-T-foc-I}:
\begin{equation}
\frac{\mu_t}{\lambda_t} = \frac{\partial \Hfun / \partial I_t}{\partial \Hfun / \partial C_t} \cdot \frac{\beta^t u'(C_t)}{\beta^t u'(C_t)} = \frac{1}{Q_t}.
\label{eq:ac-T-mu-lambda}
\end{equation}

The ratio $\mu_t / \lambda_t = 1/Q_t$ is the shadow price of investment relative to the shadow price of output. Equivalently, $\mu_t$ is the shadow value of capital (in utils), and $\lambda_t Q_t$ is the shadow value of consumption.

\paragraph{The Euler Equation (Again).}
From \cref{eq:ac-T-foc-k}: $\mu_t = \lambda_{t+1} A_{t+1}$. Substituting $\mu_t = \lambda_t / Q_t$:
\begin{equation}
\frac{\lambda_t}{Q_t} = \lambda_{t+1} A_{t+1}.
\label{eq:ac-T-lambda-recursion}
\end{equation}

From \cref{eq:ac-T-foc-C}: $\lambda_t = \beta^t u'(C_t) / (\partial \Hfun / \partial C_t)$. Using $\partial \Hfun / \partial C = 1/Q$ (up to scaling):
\begin{equation}
\lambda_t \propto \beta^t u'(C_t) Q_t.
\label{eq:ac-T-lambda-prop}
\end{equation}

Substituting into \cref{eq:ac-T-lambda-recursion} and simplifying yields the Euler equation \cref{eq:ac-T-euler}.

% --------------------------------------------------------------
\subsection{Homothetic Structure and the Normalized Problem}
\label{subsec:ac-T-homothetic}
% --------------------------------------------------------------

With homothetic preferences (CRRA) and a homogeneous-of-degree-one production/transformation technology, the problem admits a normalized formulation.

\paragraph{Guess for the Value Function.}
Conjecture that the value function takes the form:
\begin{equation}
V_t(k) = v_t \cdot \frac{k^{1-1/\theta}}{1 - 1/\theta},
\label{eq:ac-T-value-guess}
\end{equation}
where $v_t$ is a time-dependent coefficient to be determined.

\paragraph{Normalized Variables.}
Define the consumption-output ratio $\bar{c}_t \equiv C_t / Y_t$ and the investment-output ratio $\bar{\imath}_t \equiv I_t / Y_t$. The CES constraint becomes:
\begin{equation}
\Hfun(\bar{c}_t, \bar{\imath}_t) = 1,
\label{eq:ac-T-normalized-ces}
\end{equation}
which implicitly defines $\bar{c}_t$ as a function of $\bar{\imath}_t$.

\paragraph{Tobin's $Q$ in Normalized Form.}
\begin{equation}
Q_t = \frac{\omega_I^{1-\psi}}{\omega_C^{1-\psi}} \left( \frac{\bar{c}_t}{\bar{\imath}_t} \right)^{1-\psi}.
\label{eq:ac-T-Q-normalized}
\end{equation}

\paragraph{The Normalized Bellman Equation.}
Substituting the guess into the Bellman equation and using $k_{t+1} = I_t = \bar{\imath}_t Y_t = \bar{\imath}_t A_t k_t$:
\begin{equation}
v_t \frac{k^{1-1/\theta}}{1-1/\theta} = \max_{\bar{\imath}} \left\{ \frac{(\bar{c} \cdot A_t k)^{1-1/\theta}}{1-1/\theta} + \beta v_{t+1} \frac{(\bar{\imath} \cdot A_t k)^{1-1/\theta}}{1-1/\theta} \right\}.
\label{eq:ac-T-bellman-sub}
\end{equation}

Factoring out $k^{1-1/\theta}$ and $A_t^{1-1/\theta}$:
\begin{equation}
v_t = A_t^{1-1/\theta} \max_{\bar{\imath}} \left\{ \frac{\bar{c}^{1-1/\theta} + \beta v_{t+1} \bar{\imath}^{1-1/\theta}}{1-1/\theta} \right\} \cdot (1-1/\theta).
\label{eq:ac-T-bellman-normalized}
\end{equation}

This confirms the guess: the value function is linear in $k^{1-1/\theta}$, and the coefficients $v_t$ satisfy a recursion that depends only on $A_t$ and $\beta$, not on $k$.

\paragraph{The Normalized Problem.}
The coefficients $\{v_t\}$ solve a finite-dimensional problem:
\begin{equation}
v_t = \max_{\bar{\imath}_t} \left\{ \bar{c}_t^{1-1/\theta} + \beta v_{t+1} (\bar{\imath}_t A_t)^{1-1/\theta} \right\} \quad \text{s.t.} \quad \Hfun(\bar{c}_t, \bar{\imath}_t) = 1,
\label{eq:ac-T-v-recursion}
\end{equation}
with terminal condition $v_T = A_T^{1-1/\theta}$.

\begin{calloutnote}[Dimensionality Reduction]
The homothetic structure reduces the state space from $(k_t, t)$ to just $t$. The normalized investment rate $\bar{\imath}_t$ depends only on time (through $A_t$, $\beta$, and the terminal condition), not on the level of capital. This is the dynamic analog of the ``linearity in wealth'' property in consumption-savings problems.

This reduction is valuable computationally: instead of solving a value function over a continuous state space, we solve a sequence of static optimization problems, one for each period.
\end{calloutnote}

% --------------------------------------------------------------
\subsection{Steady State and Dynamics}
\label{subsec:ac-T-dynamics}
% --------------------------------------------------------------

Consider the case where productivity is constant: $A_t = A$ for all $t$. In an infinite-horizon version (or for $T$ large), the economy converges to a steady state.

\paragraph{Steady-State Conditions.}
In steady state, $\bar{c}_t = \bar{c}$, $\bar{\imath}_t = \bar{\imath}$, $Q_t = Q$, and $v_t = v$ for all $t$. The Euler equation \cref{eq:ac-T-euler} becomes:
\begin{equation}
1 = \beta \frac{A}{Q}.
\label{eq:ac-T-ss-euler}
\end{equation}

Thus, in steady state:
\begin{equation}
Q^{ss} = \beta A.
\label{eq:ac-T-Q-ss}
\end{equation}

Tobin's $Q$ equals the discounted marginal product of capital. If $\beta A > 1$, steady-state $Q > 1$: investment goods are expensive relative to the frictionless benchmark.

\paragraph{Steady-State Investment Rate.}
From the definition of $Q$:
\begin{equation}
\beta A = \frac{\omega_I^{1-\psi}}{\omega_C^{1-\psi}} \left( \frac{\bar{c}^{ss}}{\bar{\imath}^{ss}} \right)^{1-\psi}.
\label{eq:ac-T-Q-ss-def}
\end{equation}

Combined with the CES constraint $\Hfun(\bar{c}^{ss}, \bar{\imath}^{ss}) = 1$, this pins down the steady-state allocation.

\paragraph{Transitional Dynamics.}
Away from steady state, the system exhibits transitional dynamics. Starting from an initial capital stock $k_0 \neq k^{ss}$:
\begin{itemize}[nosep]
\item If $k_0 < k^{ss}$: The economy invests heavily, driving up $Q$. High $Q$ means low returns $A/Q$, which moderates investment over time.
\item If $k_0 > k^{ss}$: The economy disinvests (if possible) or lets capital depreciate, with low investment rates and low $Q$.
\end{itemize}

The adjustment cost (finite $\eta$) slows the transition: the economy cannot jump immediately to steady state because high investment rates make investment expensive at the margin.

% --------------------------------------------------------------
\subsection{Summary}
\label{subsec:ac-T-summary}
% --------------------------------------------------------------

\begin{table}[H]
\centering
\caption{The Finite-Horizon Problem with CES Adjustment Costs}
\label{tab:ac-T-summary}
\renewcommand{\arraystretch}{1.6}
\begin{tabular}{@{}ll@{}}
\toprule
\textbf{Object} & \textbf{Formula} \\
\midrule
Bellman equation & $V_t(k) = \max_{I} \left\{ u(C) + \beta V_{t+1}(I) \right\}$ s.t. $\Hfun(C,I) = A_t k$ \\[6pt]
Euler equation & $u'(C_t) = \beta \dfrac{A_{t+1}}{Q_{t+1}} Q_t \, u'(C_{t+1})$ \\[10pt]
Return on capital & $R_{t+1} = \dfrac{A_{t+1}}{Q_{t+1}} Q_t$ \\[10pt]
Steady-state $Q$ & $Q^{ss} = \beta A$ \\[6pt]
Value function form & $V_t(k) = v_t \cdot \dfrac{k^{1-1/\theta}}{1-1/\theta}$ \\[8pt]
\bottomrule
\end{tabular}
\end{table}

\begin{callouttip}[Key Insights]
The finite-horizon extension reveals:
\begin{itemize}[nosep]
\item \textbf{Returns include capital gains}: The return $R_{t+1} = (A_{t+1}/Q_{t+1}) Q_t$ has both a dividend yield and a capital gain component from changes in $Q$.
\item \textbf{Homothetic reduction}: With CRRA utility and CES transformation, the value function is homogeneous in capital, reducing the problem to a sequence of normalized static problems.
\item \textbf{Steady-state $Q$}: In steady state, $Q = \beta A$---Tobin's $Q$ equals the discounted marginal product.
\item \textbf{Gradual adjustment}: Finite transformation elasticity $\eta$ implies gradual convergence to steady state; the economy cannot jump instantaneously.
\end{itemize}
\end{callouttip}


% ==============================================================
\section{Capital Heterogeneity}
\label{sec:ac-heterogeneity}
% ==============================================================

We now introduce multiple capital types, indexed by $i \in \{1, 2, \ldots, N\}$. This extension connects adjustment costs to portfolio theory: the planner allocates investment across sectors, facing a trade-off between diversification and the curvature of the production technology. The CES structure appears at two levels---in the transformation of output into consumption and investment, and in the aggregation of different capital types into output.

% --------------------------------------------------------------
\subsection{Setup with Multiple Capital Types}
\label{subsec:ac-het-setup}
% --------------------------------------------------------------

\paragraph{Technology.}
The economy has $N$ types of capital, with stocks $\bk = (k_1, \ldots, k_N)$. Output is produced by a CES aggregator over effective capital inputs:
\begin{equation}
Y = \F(\bk) = \left( \sum_{i=1}^{N} \alpha_i^{1-\sigma} k_i^{\sigma} \right)^{1/\sigma},
\label{eq:ac-het-production}
\end{equation}
where $\sigma \in (-\infty, 1]$ is the power parameter and $\alpha_i > 0$ are weights satisfying $\sum_i \alpha_i = 1$. The elasticity of substitution across capital types is:
\begin{equation}
\varepsilon = \frac{1}{1 - \sigma}.
\label{eq:ac-het-epsilon}
\end{equation}

\paragraph{Transformation Technology.}
Output is transformed into consumption $C$ and a vector of investments $\bI = (I_1, \ldots, I_N)$ through a CES aggregator:
\begin{equation}
\Hfun(C, \bI) = \left( \omega_C^{1-\psi} C^{\psi} + \sum_{i=1}^{N} \omega_i^{1-\psi} I_i^{\psi} \right)^{1/\psi} = Y,
\label{eq:ac-het-transformation}
\end{equation}
where $\psi$ governs the transformation elasticity $\eta = 1/(1-\psi)$, and the weights satisfy $\omega_C + \sum_i \omega_i = 1$.

\paragraph{Capital Accumulation.}
With full depreciation ($\delta = 1$):
\begin{equation}
k_i' = I_i, \qquad i = 1, \ldots, N.
\label{eq:ac-het-accum}
\end{equation}

\paragraph{Notation.}
We use bold for vectors: $\bk = (k_1, \ldots, k_N)$, $\bI = (I_1, \ldots, I_N)$, $\boldsymbol{\alpha} = (\alpha_1, \ldots, \alpha_N)$, $\boldsymbol{\omega} = (\omega_1, \ldots, \omega_N)$. Let $\be$ denote the vector of ones, so aggregate capital is $K = \be^\top \bk = \sum_i k_i$.

% --------------------------------------------------------------
\subsection{The Two-Period Problem}
\label{subsec:ac-het-two-period}
% --------------------------------------------------------------

Consider the two-period problem. The planner enters period 0 with capital stocks $\bk_0$ and chooses $(C_0, \bI_0)$ to maximize:
\begin{equation}
\max_{C_0, \bI_0} \left\{ u(C_0) + \beta \, u(C_1) \right\}
\label{eq:ac-het-objective}
\end{equation}
subject to:
\begin{align}
\Hfun(C_0, \bI_0) &= \F(\bk_0), \label{eq:ac-het-constraint0} \\
C_1 &= \F(\bI_0). \label{eq:ac-het-constraint1}
\end{align}

\paragraph{The Nested Structure.}
The problem has a nested structure:
\begin{enumerate}[nosep]
\item \textbf{Outer problem}: Allocate output $Y_0$ between consumption $C_0$ and total investment.
\item \textbf{Inner problem}: Allocate total investment across capital types $\bI_0$ to maximize future output $\F(\bI_0)$.
\end{enumerate}

This decomposition is possible because future consumption $C_1 = \F(\bI_0)$ depends only on the composition of investment, not on the consumption-investment split.

% --------------------------------------------------------------
\subsection{The Portfolio Problem}
\label{subsec:ac-het-portfolio}
% --------------------------------------------------------------

\paragraph{Defining Portfolio Shares.}
Let $\bar{I} = \sum_i I_i$ denote total investment and define portfolio shares:
\begin{equation}
\phi_i \equiv \frac{I_i}{\bar{I}}, \qquad \sum_{i=1}^{N} \phi_i = 1.
\label{eq:ac-het-shares}
\end{equation}

Future output can be written as:
\begin{equation}
\F(\bI) = \F(\phi_1 \bar{I}, \ldots, \phi_N \bar{I}) = \bar{I} \cdot \F(\boldsymbol{\phi}),
\label{eq:ac-het-F-homog}
\end{equation}
using the homogeneity of $\F$. Here $\F(\boldsymbol{\phi})$ is the ``return per unit of investment'' given portfolio $\boldsymbol{\phi}$:
\begin{equation}
\F(\boldsymbol{\phi}) = \left( \sum_{i=1}^{N} \alpha_i^{1-\sigma} \phi_i^{\sigma} \right)^{1/\sigma}.
\label{eq:ac-het-F-phi}
\end{equation}

\paragraph{The Inner Problem: Optimal Portfolio.}
The planner chooses $\boldsymbol{\phi}$ to maximize the return per unit of investment:
\begin{equation}
\max_{\boldsymbol{\phi}} \; \F(\boldsymbol{\phi}) = \left( \sum_{i=1}^{N} \alpha_i^{1-\sigma} \phi_i^{\sigma} \right)^{1/\sigma} \quad \text{subject to} \quad \sum_{i=1}^{N} \phi_i = 1.
\label{eq:ac-het-inner}
\end{equation}

This is exactly the canonical CES optimization problem! Applying our standard results:

\begin{proposition}[Optimal Investment Portfolio]
\label{prop:ac-het-portfolio}
The optimal portfolio shares are:
\begin{equation}
\phi_i^* = \alpha_i.
\label{eq:ac-het-phi-star}
\end{equation}
The maximized return per unit of investment is:
\begin{equation}
\F^* \equiv \F(\boldsymbol{\alpha}) = \left( \sum_{i=1}^{N} \alpha_i^{1-\sigma} \alpha_i^{\sigma} \right)^{1/\sigma} = \left( \sum_{i=1}^{N} \alpha_i \right)^{1/\sigma} = 1.
\label{eq:ac-het-F-star}
\end{equation}
\end{proposition}

\begin{proof}
This follows directly from the benchmark solution (Proposition~\ref{prop:benchmark}): with equal ``prices'' (the constraint $\sum_i \phi_i = 1$ has uniform coefficients), the optimal allocation equals the canonical weights.
\end{proof}

\begin{calloutnote}[Portfolio Separation]
The optimal portfolio $\phi_i^* = \alpha_i$ is \emph{independent} of:
\begin{itemize}[nosep]
\item The level of total investment $\bar{I}$.
\item The transformation elasticity $\eta$.
\item The discount factor $\beta$ and IES $\theta$.
\end{itemize}
This is a separation result: the composition of investment is determined solely by the production technology (through $\boldsymbol{\alpha}$), while the level of investment is determined by the consumption-savings trade-off.

The result $\F^* = 1$ (with our normalization $\sum_i \alpha_i = 1$) means that one unit of optimally allocated investment produces one unit of future output.
\end{calloutnote}

% --------------------------------------------------------------
\subsection{The Outer Problem: Consumption vs. Investment}
\label{subsec:ac-het-outer}
% --------------------------------------------------------------

Given the optimal portfolio, future consumption is $C_1 = \F(\bI_0) = \bar{I}_0 \cdot \F^* = \bar{I}_0$ (using $\F^* = 1$). The problem reduces to choosing $(C_0, \bar{I}_0)$:
\begin{equation}
\max_{C_0, \bar{I}_0} \left\{ u(C_0) + \beta \, u(\bar{I}_0) \right\} \quad \text{subject to} \quad \Hfun(C_0, \bI_0) = Y_0,
\label{eq:ac-het-outer}
\end{equation}
where $\bI_0 = \boldsymbol{\alpha} \bar{I}_0$ (the optimal portfolio).

\paragraph{The Effective Transformation.}
Substituting $I_i = \alpha_i \bar{I}$ into the transformation function:
\begin{equation}
\Hfun(C, \boldsymbol{\alpha} \bar{I}) = \left( \omega_C^{1-\psi} C^{\psi} + \sum_{i=1}^{N} \omega_i^{1-\psi} (\alpha_i \bar{I})^{\psi} \right)^{1/\psi}.
\label{eq:ac-het-H-sub}
\end{equation}

Define the \emph{effective investment weight}:
\begin{equation}
\tilde{\omega}_I \equiv \left( \sum_{i=1}^{N} \omega_i^{1-\psi} \alpha_i^{\psi} \right)^{1/(1-\psi)}.
\label{eq:ac-het-omega-tilde}
\end{equation}

Then the transformation simplifies to:
\begin{equation}
\Hfun(C, \boldsymbol{\alpha} \bar{I}) = \left( \omega_C^{1-\psi} C^{\psi} + \tilde{\omega}_I^{1-\psi} \bar{I}^{\psi} \right)^{1/\psi} \equiv \tilde{\Hfun}(C, \bar{I}).
\label{eq:ac-het-H-effective}
\end{equation}

\begin{calloutimportant}[Aggregation Result]
When investment is optimally allocated, the multi-sector transformation technology aggregates to an \emph{equivalent two-good technology} $\tilde{\Hfun}(C, \bar{I})$ with effective weight $\tilde{\omega}_I$.

The effective weight $\tilde{\omega}_I$ combines the transformation weights $\omega_i$ and the production weights $\alpha_i$. It is itself a CES aggregate---a generalized mean of the $\omega_i$ with weights $\alpha_i$:
\[
\tilde{\omega}_I = \left( \sum_{i=1}^{N} \alpha_i \left( \frac{\omega_i}{\alpha_i} \right)^{1-\psi} \alpha_i^{\psi} \right)^{1/(1-\psi)} = \M\left( \frac{\boldsymbol{\omega}}{\boldsymbol{\alpha}}; \boldsymbol{\alpha}, \psi \right) \cdot \left( \sum_i \alpha_i \right).
\]
This nesting of CES aggregators is a recurring theme.
\end{calloutimportant}

\paragraph{Tobin's $Q$ with Heterogeneous Capital.}
With the effective transformation $\tilde{\Hfun}(C, \bar{I}) = Y$, Tobin's $Q$ for aggregate investment is:
\begin{equation}
Q = \frac{\tilde{\omega}_I^{1-\psi}}{\omega_C^{1-\psi}} \left( \frac{C}{\bar{I}} \right)^{1-\psi}.
\label{eq:ac-het-Q}
\end{equation}

The Euler equation takes the same form as in the single-capital case:
\begin{equation}
u'(C_0) = \beta \frac{1}{Q} u'(C_1),
\label{eq:ac-het-euler}
\end{equation}
where we used $\F^* = 1$ so that the ``marginal product'' of optimally allocated investment is unity.

% --------------------------------------------------------------
\subsection{Sector-Specific Prices}
\label{subsec:ac-het-prices}
% --------------------------------------------------------------

\paragraph{Decentralization.}
To decentralize, we need prices for each investment good $I_i$. The transformation firm solves:
\begin{equation}
\max_{C, \bI} \; p_C C + \sum_{i=1}^{N} p_i I_i \quad \text{subject to} \quad \Hfun(C, \bI) = Y.
\label{eq:ac-het-firm}
\end{equation}

The first-order conditions give:
\begin{equation}
\frac{p_i}{p_C} = \frac{\partial \Hfun / \partial I_i}{\partial \Hfun / \partial C} = \frac{\omega_i^{1-\psi}}{\omega_C^{1-\psi}} \left( \frac{C}{I_i} \right)^{1-\psi}.
\label{eq:ac-het-price-ratio}
\end{equation}

\paragraph{Sector-Specific Tobin's $Q$.}
Define the sector-specific Tobin's $Q$:
\begin{equation}
Q_i \equiv \frac{p_i}{p_C} = \frac{\omega_i^{1-\psi}}{\omega_C^{1-\psi}} \left( \frac{C}{I_i} \right)^{1-\psi}.
\label{eq:ac-het-Qi}
\end{equation}

At the optimal portfolio where $I_i = \alpha_i \bar{I}$:
\begin{equation}
Q_i^* = \frac{\omega_i^{1-\psi}}{\omega_C^{1-\psi}} \left( \frac{C}{\alpha_i \bar{I}} \right)^{1-\psi} = \frac{\omega_i^{1-\psi}}{\alpha_i^{1-\psi}} \cdot \frac{1}{\omega_C^{1-\psi}} \left( \frac{C}{\bar{I}} \right)^{1-\psi}.
\label{eq:ac-het-Qi-star}
\end{equation}

\paragraph{Relative Prices Across Sectors.}
The ratio of sector prices is:
\begin{equation}
\frac{Q_i}{Q_j} = \frac{\omega_i^{1-\psi} / \alpha_i^{1-\psi}}{\omega_j^{1-\psi} / \alpha_j^{1-\psi}} = \left( \frac{\omega_i / \alpha_i}{\omega_j / \alpha_j} \right)^{1-\psi}.
\label{eq:ac-het-price-ratio-ij}
\end{equation}

When $\omega_i / \alpha_i = \omega_j / \alpha_j$ for all $i, j$ (i.e., transformation weights are proportional to production weights), all investment goods have the same price: $Q_i = Q_j = Q$.

\begin{calloutnote}[When Prices Differ]
Sector-specific prices $Q_i$ differ when the transformation weights $\omega_i$ are not proportional to the production weights $\alpha_i$. Intuitively:
\begin{itemize}[nosep]
\item If $\omega_i / \alpha_i$ is high, sector $i$ is ``easy to produce'' relative to its importance in production. Its price $Q_i$ is low.
\item If $\omega_i / \alpha_i$ is low, sector $i$ is ``hard to produce.'' Its price $Q_i$ is high.
\end{itemize}

In the symmetric case ($\omega_i = \alpha_i$ for all $i$), all sectors have equal prices and the multi-sector model reduces exactly to the single-sector case.
\end{calloutnote}

% --------------------------------------------------------------
\subsection{The Production Firm's Problem}
\label{subsec:ac-het-production-firm}
% --------------------------------------------------------------

A separate production firm (or the same firm in a later period) uses installed capital to produce output. The firm solves:
\begin{equation}
\max_{\bk} \; p_Y \F(\bk) - \sum_{i=1}^{N} r_i k_i,
\label{eq:ac-het-prod-firm}
\end{equation}
where $r_i$ is the rental rate for capital of type $i$.

\paragraph{First-Order Conditions.}
\begin{equation}
r_i = p_Y \frac{\partial \F}{\partial k_i} = p_Y \alpha_i^{1-\sigma} k_i^{\sigma - 1} \F^{1-\sigma}.
\label{eq:ac-het-rental}
\end{equation}

\paragraph{Zero Profits.}
By homogeneity of $\F$:
\begin{equation}
p_Y \F(\bk) = \sum_{i=1}^{N} r_i k_i.
\label{eq:ac-het-prod-zero-profit}
\end{equation}

The production firm earns zero profits, as expected with CRS technology.

\paragraph{Equilibrium Rental Rates.}
With full depreciation, $k_i = I_{i,-1}$ (capital equals last period's investment). The rental rate satisfies:
\begin{equation}
r_i = p_Y \alpha_i^{1-\sigma} \left( \frac{k_i}{\F(\bk)} \right)^{\sigma - 1} = p_Y \alpha_i^{1-\sigma} \left( \frac{I_{i,-1}}{Y} \right)^{\sigma - 1}.
\label{eq:ac-het-rental-eq}
\end{equation}

At the optimal portfolio where $I_i / \bar{I} = \alpha_i$:
\begin{equation}
r_i^* = p_Y \alpha_i^{1-\sigma} \alpha_i^{\sigma - 1} = p_Y.
\label{eq:ac-het-rental-star}
\end{equation}

All capital types earn the same rental rate $r_i = p_Y$ when investment is optimally allocated!

\begin{calloutimportant}[Equalization of Rental Rates]
At the optimal portfolio, the marginal product of each capital type (weighted by its price) is equalized:
\[
\frac{r_i}{p_Y} = \frac{\partial \F}{\partial k_i} = 1 \quad \text{for all } i.
\]
This is the multi-sector analog of the single-sector condition $\partial F / \partial k = A$. Optimal diversification implies that no sector offers a higher risk-adjusted return than any other.
\end{calloutimportant}

% --------------------------------------------------------------
\subsection{Connection to Portfolio Theory}
\label{subsec:ac-het-portfolio-theory}
% --------------------------------------------------------------

The structure of the investment allocation problem mirrors portfolio choice in finance.

\paragraph{The Analogy.}
\begin{center}
\begin{tabular}{@{}lll@{}}
\toprule
\textbf{Concept} & \textbf{Portfolio Theory} & \textbf{Capital Allocation} \\
\midrule
Choice variable & Asset shares $\alpha_j$ & Investment shares $\phi_i$ \\
Constraint & $\sum_j \alpha_j = 1$ & $\sum_i \phi_i = 1$ \\
Objective & Certainty equivalent of returns & Production $\F(\boldsymbol{\phi})$ \\
Aggregator & $\CE(\bR; \boldsymbol{\pi}, \gamma)$ & $\F(\boldsymbol{\phi}; \boldsymbol{\alpha}, \sigma)$ \\
Curvature & Risk aversion $\gamma$ & Substitution $\varepsilon = 1/(1-\sigma)$ \\
\bottomrule
\end{tabular}
\end{center}

\paragraph{Key Difference: Deterministic vs. Stochastic.}
In portfolio theory, the curvature of the certainty equivalent reflects risk aversion over uncertain outcomes. Here, the curvature of $\F$ reflects technological complementarity/substitutability across capital types---a deterministic feature.

However, when we introduce uncertainty (Section~\ref{sec:ac-risk}), the two interpretations merge: the production function aggregates over states (via capital-specific productivity shocks), and the curvature reflects both technological and risk considerations.

\paragraph{The CES as a Unifying Framework.}
The CES aggregator provides a unified treatment:
\begin{itemize}[nosep]
\item \textbf{Preferences}: Elasticity of substitution across goods, time, or states.
\item \textbf{Technology}: Elasticity of substitution across inputs.
\item \textbf{Transformation}: Elasticity of transformation between consumption and investment.
\end{itemize}

The optimal allocation rule $\phi_i^* = \alpha_i$ (invest proportionally to production weights) is the deterministic analog of ``invest proportionally to probabilities'' in the log-utility portfolio problem.

% --------------------------------------------------------------
\subsection{Dynamics with Capital Heterogeneity}
\label{subsec:ac-het-dynamics}
% --------------------------------------------------------------

\paragraph{The $T$-Period Problem.}
Extending to $T$ periods, the planner solves:
\begin{equation}
\max_{\{C_t, \bI_t\}} \sum_{t=0}^{T} \beta^t u(C_t)
\label{eq:ac-het-T-obj}
\end{equation}
subject to:
\begin{align}
\Hfun(C_t, \bI_t) &= \F(\bk_t), \qquad t = 0, \ldots, T-1, \\
\bk_{t+1} &= \bI_t, \\
C_T &= \F(\bk_T).
\end{align}

\paragraph{Portfolio Separation Over Time.}
At each date $t$, the optimal investment portfolio is $\phi_{i,t}^* = \alpha_i$---independent of $t$. The portfolio is constant over time, even as the level of investment $\bar{I}_t$ varies.

This follows because the production function $\F$ is time-invariant (no capital-specific productivity shocks). The optimal portfolio maximizes the static return $\F(\boldsymbol{\phi})$, which has the same solution at every date.

\paragraph{Aggregate Dynamics.}
Define aggregate investment $\bar{I}_t = \sum_i I_{i,t}$. The aggregate dynamics follow:
\begin{align}
\tilde{\Hfun}(C_t, \bar{I}_t) &= Y_t, \\
Y_{t+1} &= \F(\boldsymbol{\alpha} \bar{I}_t) = \bar{I}_t,
\end{align}
where we used $\F(\boldsymbol{\alpha}) = 1$.

This is exactly the single-capital model with $A = 1$! The heterogeneous-capital economy, under optimal allocation, behaves like a single-capital economy with unit productivity.

\begin{callouttip}[Aggregation Theorem]
With CES production and transformation technologies, the heterogeneous-capital economy aggregates to an equivalent single-capital economy:
\begin{itemize}[nosep]
\item The optimal portfolio is $\phi_i^* = \alpha_i$ (time-invariant).
\item The effective transformation is $\tilde{\Hfun}(C, \bar{I}) = Y$ with weight $\tilde{\omega}_I$.
\item The effective productivity is $\F^* = 1$ (with normalized weights).
\item Aggregate dynamics follow the single-capital model of Section~\ref{sec:ac-T-period}.
\end{itemize}
This aggregation result relies on homogeneity and the absence of capital-specific shocks. When we introduce risk (next section), the aggregation becomes more subtle.
\end{callouttip}

% --------------------------------------------------------------
\subsection{Summary}
\label{subsec:ac-het-summary}
% --------------------------------------------------------------

\begin{table}[H]
\centering
\caption{Capital Heterogeneity with CES Adjustment Costs}
\label{tab:ac-het-summary}
\renewcommand{\arraystretch}{1.6}
\begin{tabular}{@{}ll@{}}
\toprule
\textbf{Object} & \textbf{Formula} \\
\midrule
Production & $Y = \F(\bk) = \left( \sum_i \alpha_i^{1-\sigma} k_i^{\sigma} \right)^{1/\sigma}$ \\[6pt]
Transformation & $\Hfun(C, \bI) = \left( \omega_C^{1-\psi} C^{\psi} + \sum_i \omega_i^{1-\psi} I_i^{\psi} \right)^{1/\psi} = Y$ \\[6pt]
Optimal portfolio & $\phi_i^* = I_i / \bar{I} = \alpha_i$ \\[6pt]
Effective weight & $\tilde{\omega}_I = \left( \sum_i \omega_i^{1-\psi} \alpha_i^{\psi} \right)^{1/(1-\psi)}$ \\[6pt]
Sector price & $Q_i = \dfrac{\omega_i^{1-\psi}}{\omega_C^{1-\psi}} \left( \dfrac{C}{I_i} \right)^{1-\psi}$ \\[10pt]
Aggregate $Q$ & $Q = \dfrac{\tilde{\omega}_I^{1-\psi}}{\omega_C^{1-\psi}} \left( \dfrac{C}{\bar{I}} \right)^{1-\psi}$ \\[10pt]
\bottomrule
\end{tabular}
\end{table}

\begin{callouttip}[Key Insights]
Capital heterogeneity with CES technologies yields:
\begin{itemize}[nosep]
\item \textbf{Portfolio separation}: The optimal investment mix $\phi_i^* = \alpha_i$ is independent of the consumption-savings decision and constant over time.
\item \textbf{Aggregation}: The multi-sector economy reduces to an equivalent single-sector economy with effective weight $\tilde{\omega}_I$.
\item \textbf{Sector-specific prices}: Each capital type has its own Tobin's $Q_i$, which differ when $\omega_i / \alpha_i$ varies across sectors.
\item \textbf{Zero profits}: Both transformation and production firms earn zero profits under CRS.
\item \textbf{Connection to portfolio theory}: The investment allocation problem has the same structure as portfolio choice, with production curvature playing the role of risk aversion.
\end{itemize}
\end{callouttip}


% ==============================================================
\section{Incomplete Depreciation}
\label{sec:ac-depreciation}
% ==============================================================

We now relax the assumption of full depreciation. When capital persists across periods ($\delta < 1$), the distinction between gross and net investment becomes important, and the dynamics of Tobin's $Q$ acquire a richer structure. The CES transformation technology continues to govern adjustment costs, but now applies to gross investment rather than the change in the capital stock.

% --------------------------------------------------------------
\subsection{Setup with Persistent Capital}
\label{subsec:ac-dep-setup}
% --------------------------------------------------------------

\paragraph{Capital Accumulation.}
Capital depreciates at rate $\delta \in (0, 1]$ and is augmented by gross investment $I$:
\begin{equation}
k' = (1 - \delta) k + I.
\label{eq:ac-dep-accum}
\end{equation}

Net investment is the change in the capital stock:
\begin{equation}
k' - k = I - \delta k.
\label{eq:ac-dep-net}
\end{equation}

When $\delta = 1$, we recover the full depreciation case: $k' = I$.

\paragraph{Technology.}
Production is linear in capital:
\begin{equation}
Y = A k.
\label{eq:ac-dep-production}
\end{equation}

The CES transformation governs the split of output into consumption and gross investment:
\begin{equation}
\Hfun(C, I) = \left( \omega_C^{1-\psi} C^{\psi} + \omega_I^{1-\psi} I^{\psi} \right)^{1/\psi} = Y.
\label{eq:ac-dep-transformation}
\end{equation}

\paragraph{Timing.}
At the beginning of period $t$, the household owns capital $k_t$. Production occurs, yielding output $Y_t = A k_t$. The household then chooses consumption $C_t$ and gross investment $I_t$ subject to the transformation constraint. Capital for next period is $k_{t+1} = (1-\delta) k_t + I_t$.

% --------------------------------------------------------------
\subsection{The Two-Period Problem}
\label{subsec:ac-dep-two-period}
% --------------------------------------------------------------

\paragraph{Setup.}
Consider $t \in \{0, 1\}$. The household enters period 0 with capital $k_0$, chooses $(C_0, I_0)$, and then consumes all output in period 1:
\begin{equation}
\max_{C_0, I_0} \left\{ u(C_0) + \beta \, u(C_1) \right\}
\label{eq:ac-dep-objective}
\end{equation}
subject to:
\begin{align}
\Hfun(C_0, I_0) &= A k_0, \label{eq:ac-dep-const0} \\
k_1 &= (1-\delta) k_0 + I_0, \label{eq:ac-dep-accum0} \\
C_1 &= A k_1 = A \left[ (1-\delta) k_0 + I_0 \right]. \label{eq:ac-dep-const1}
\end{align}

\paragraph{The Key Difference.}
With incomplete depreciation, future consumption depends on \emph{both} current investment $I_0$ and the undepreciated capital $(1-\delta) k_0$:
\begin{equation}
C_1 = A (1-\delta) k_0 + A I_0.
\label{eq:ac-dep-C1}
\end{equation}

The term $A(1-\delta) k_0$ is ``locked in''---it does not depend on the household's choice. Only the term $A I_0$ is affected by the investment decision.

\paragraph{First-Order Condition.}
The Lagrangian is:
\begin{equation}
\mathcal{L} = u(C_0) + \beta \, u(A[(1-\delta)k_0 + I_0]) + \lambda \left[ A k_0 - \Hfun(C_0, I_0) \right].
\label{eq:ac-dep-lagrangian}
\end{equation}

First-order conditions:
\begin{align}
C_0: & \quad u'(C_0) = \lambda \frac{\partial \Hfun}{\partial C_0}, \label{eq:ac-dep-foc-C} \\
I_0: & \quad \beta A \, u'(C_1) = \lambda \frac{\partial \Hfun}{\partial I_0}. \label{eq:ac-dep-foc-I}
\end{align}

Taking the ratio and using the definition of Tobin's $Q$:
\begin{equation}
\frac{u'(C_0)}{\beta A \, u'(C_1)} = \frac{\partial \Hfun / \partial C_0}{\partial \Hfun / \partial I_0} = \frac{1}{Q}.
\label{eq:ac-dep-ratio}
\end{equation}

The Euler equation is:
\begin{equation}
\boxed{u'(C_0) = \beta \frac{A}{Q} u'(C_1).}
\label{eq:ac-dep-euler}
\end{equation}

This has the same form as the full depreciation case! The Euler equation relates marginal utilities through the return $A/Q$, regardless of $\delta$.

\begin{calloutnote}[Why Doesn't $\delta$ Appear in the Euler Equation?]
The depreciation rate $\delta$ does not appear in the Euler equation because it affects only the \emph{level} of future consumption, not the \emph{marginal} trade-off. The household's decision at the margin is: ``Should I invest one more unit today?'' The return on that marginal unit is $A/Q$, independent of how much capital survives from the past.

However, $\delta$ does affect the \emph{solution} through the budget constraint: higher $\delta$ means less future consumption for given $I_0$, which changes the equilibrium allocation and hence $Q$.
\end{calloutnote}

% --------------------------------------------------------------
\subsection{The Finite-Horizon Problem}
\label{subsec:ac-dep-T-period}
% --------------------------------------------------------------

\paragraph{The Bellman Equation.}
For $t \in \{0, 1, \ldots, T\}$, define $V_t(k)$ as the value of entering period $t$ with capital $k$:
\begin{equation}
V_t(k) = \max_{I} \left\{ u(C) + \beta V_{t+1}((1-\delta)k + I) \right\} \quad \text{s.t.} \quad \Hfun(C, I) = A k,
\label{eq:ac-dep-bellman}
\end{equation}
with terminal condition $V_T(k) = u(Ak)$.

\paragraph{First-Order Condition.}
\begin{equation}
-u'(C) \frac{1}{Q} + \beta V_{t+1}'(k') = 0,
\label{eq:ac-dep-foc-bellman}
\end{equation}
where $k' = (1-\delta)k + I$.

\paragraph{Envelope Condition.}
Differentiating the value function:
\begin{equation}
V_t'(k) = u'(C) \frac{A}{Q} + \beta (1-\delta) V_{t+1}'(k').
\label{eq:ac-dep-envelope}
\end{equation}

The envelope condition now has two terms:
\begin{itemize}[nosep]
\item $u'(C) \cdot A/Q$: The marginal value of capital through current production (as before).
\item $\beta (1-\delta) V_{t+1}'(k')$: The marginal value of capital that survives to next period.
\end{itemize}

\paragraph{The Euler Equation.}
Combining the FOC and envelope condition (shifted forward):
\begin{equation}
\frac{u'(C_t)}{Q_t} = \beta V_{t+1}'(k_{t+1}) = \beta \left[ u'(C_{t+1}) \frac{A}{Q_{t+1}} + \beta (1-\delta) V_{t+2}'(k_{t+2}) \right].
\label{eq:ac-dep-euler-step}
\end{equation}

Iterating forward and using the FOC repeatedly:
\begin{equation}
\frac{u'(C_t)}{Q_t} = \beta u'(C_{t+1}) \frac{A}{Q_{t+1}} + \beta^2 (1-\delta) u'(C_{t+2}) \frac{A}{Q_{t+2}} + \cdots
\label{eq:ac-dep-euler-iterated}
\end{equation}

This can be written more compactly. Define the \emph{user cost of capital}:
\begin{equation}
\boxed{u'(C_t) = \beta \left[ \frac{A}{Q_{t+1}} + (1-\delta) \frac{Q_t}{Q_{t+1}} \right] Q_t \, u'(C_{t+1}).}
\label{eq:ac-dep-euler-full}
\end{equation}

\paragraph{Interpretation: The Return on Capital.}
Define the gross return on capital:
\begin{equation}
R_{t+1} \equiv \frac{A + (1-\delta) Q_{t+1}}{Q_t}.
\label{eq:ac-dep-return}
\end{equation}

The Euler equation becomes:
\begin{equation}
u'(C_t) = \beta R_{t+1} \, u'(C_{t+1}).
\label{eq:ac-dep-euler-standard}
\end{equation}

\begin{calloutimportant}[The Return with Incomplete Depreciation]
The return $R_{t+1} = (A + (1-\delta) Q_{t+1}) / Q_t$ has two components:
\begin{enumerate}[nosep]
\item \textbf{Dividend yield}: $A / Q_t$ is the marginal product per unit of capital value.
\item \textbf{Capital gain}: $(1-\delta) Q_{t+1} / Q_t$ is the value of undepreciated capital, adjusted for changes in $Q$.
\end{enumerate}

Compare with full depreciation where $R_{t+1} = A Q_t / Q_{t+1}$. With $\delta < 1$, the ``resale value'' of capital $(1-\delta) Q_{t+1}$ enters the return.

In steady state where $Q_t = Q_{t+1} = Q$:
\[
R = \frac{A + (1-\delta) Q}{Q} = \frac{A}{Q} + (1-\delta).
\]
The return equals the dividend yield plus the survival rate.
\end{calloutimportant}

% --------------------------------------------------------------
\subsection{Steady State}
\label{subsec:ac-dep-steady-state}
% --------------------------------------------------------------

Consider constant productivity $A$ and an infinite horizon (or large $T$). In steady state, all variables are constant: $C_t = C$, $I_t = I$, $k_t = k$, $Q_t = Q$.

\paragraph{Steady-State Capital Accumulation.}
From $k' = (1-\delta)k + I$:
\begin{equation}
k = (1-\delta)k + I \quad \Longrightarrow \quad I = \delta k.
\label{eq:ac-dep-ss-accum}
\end{equation}

In steady state, gross investment exactly replaces depreciated capital. Net investment is zero.

\paragraph{Steady-State Euler Equation.}
From \cref{eq:ac-dep-euler-standard} with $R_{t+1} = R$:
\begin{equation}
1 = \beta R = \beta \frac{A + (1-\delta) Q}{Q}.
\label{eq:ac-dep-ss-euler}
\end{equation}

Solving for $Q$:
\begin{equation}
Q = \frac{\beta A}{1 - \beta(1-\delta)}.
\label{eq:ac-dep-Q-ss}
\end{equation}

\paragraph{Special Cases.}

\medskip
\noindent
\textit{Full depreciation ($\delta = 1$):}
\begin{equation}
Q^{ss} = \frac{\beta A}{1 - 0} = \beta A.
\label{eq:ac-dep-Q-ss-full}
\end{equation}
This matches our earlier result.

\medskip
\noindent
\textit{No depreciation ($\delta = 0$):}
\begin{equation}
Q^{ss} = \frac{\beta A}{1 - \beta} = \frac{A}{1/\beta - 1} = \frac{A}{r},
\label{eq:ac-dep-Q-ss-none}
\end{equation}
where $r = 1/\beta - 1$ is the discount rate. This is the standard asset pricing formula: the price equals the perpetuity value of dividends.

\medskip
\noindent
\textit{General $\delta \in (0,1)$:}
\begin{equation}
Q^{ss} = \frac{A}{1/\beta - (1-\delta)} = \frac{A}{r + \delta}.
\label{eq:ac-dep-Q-ss-general}
\end{equation}

The denominator $r + \delta$ is the \emph{user cost of capital}: the sum of the interest cost and the depreciation cost.

\begin{calloutnote}[The User Cost of Capital]
The steady-state condition $Q = A / (r + \delta)$ can be rewritten as:
\[
(r + \delta) Q = A.
\]
The left-hand side is the \emph{user cost}---the cost of using one unit of capital for one period:
\begin{itemize}[nosep]
\item $r Q$: Interest cost (opportunity cost of funds tied up in capital).
\item $\delta Q$: Depreciation cost (value of capital lost to wear).
\end{itemize}
In equilibrium, the user cost equals the marginal product $A$. This is the classic Jorgensonian result, extended to include adjustment costs through $Q$.
\end{calloutnote}

\paragraph{Steady-State Investment Rate.}
From $I = \delta k$ and $Y = Ak$:
\begin{equation}
\frac{I}{Y} = \frac{\delta k}{A k} = \frac{\delta}{A}.
\label{eq:ac-dep-inv-rate}
\end{equation}

The investment-output ratio is $\delta / A$. Higher depreciation requires more investment to maintain the capital stock.

\paragraph{Steady-State Consumption.}
From the CES constraint:
\begin{equation}
\Hfun(C, I) = Y \quad \Longrightarrow \quad \Hfun(C, \delta k) = A k.
\label{eq:ac-dep-ss-ces}
\end{equation}

Given $Q$, we can solve for the consumption-investment ratio and hence for $C$ and $k$.

% --------------------------------------------------------------
\subsection{Transitional Dynamics}
\label{subsec:ac-dep-dynamics}
% --------------------------------------------------------------

Away from steady state, the system exhibits transitional dynamics governed by two state-like variables: the capital stock $k_t$ and (implicitly) the shadow value $V_t'(k_t)$.

\paragraph{The Dynamic System.}
The economy is characterized by:
\begin{align}
k_{t+1} &= (1-\delta) k_t + I_t, \label{eq:ac-dep-dyn-k} \\
\frac{u'(C_t)}{Q_t} &= \beta \left[ \frac{A}{Q_{t+1}} u'(C_{t+1}) + (1-\delta) \frac{u'(C_{t+1})}{Q_{t+1}} \right], \label{eq:ac-dep-dyn-euler} \\
\Hfun(C_t, I_t) &= A k_t. \label{eq:ac-dep-dyn-ces}
\end{align}

Given $k_t$, the CES constraint \cref{eq:ac-dep-dyn-ces} defines a relationship between $C_t$ and $I_t$. The Euler equation \cref{eq:ac-dep-dyn-euler} pins down the intertemporal allocation. Capital accumulation \cref{eq:ac-dep-dyn-k} determines $k_{t+1}$.

\paragraph{Phase Diagram Intuition.}
In the $(k, C)$ plane:
\begin{itemize}[nosep]
\item The $\dot{k} = 0$ locus (where $k_{t+1} = k_t$) requires $I = \delta k$.
\item The $\dot{C} = 0$ locus (where $C_{t+1} = C_t$) requires $\beta R = 1$, which pins down $Q$ and hence the $C/I$ ratio.
\end{itemize}

The steady state is at the intersection. The saddle-path dynamics are qualitatively similar to the standard Ramsey model, but the adjustment cost (finite $\eta$) affects the speed of convergence.

\paragraph{Effect of Adjustment Costs on Dynamics.}
Lower transformation elasticity $\eta$ (higher adjustment costs):
\begin{itemize}[nosep]
\item Slows convergence to steady state.
\item Increases the volatility of $Q$ along the transition path.
\item Reduces the response of investment to productivity shocks.
\end{itemize}

In the limit $\eta \to \infty$ (no adjustment costs), the economy jumps immediately to steady state. In the limit $\eta \to 0$ (Leontief), the consumption-investment ratio is fixed and adjustment is impossible.

% --------------------------------------------------------------
\subsection{Gross vs. Net Investment}
\label{subsec:ac-dep-gross-net}
% --------------------------------------------------------------

With incomplete depreciation, the distinction between gross and net investment matters for measurement and interpretation.

\paragraph{Definitions.}
\begin{align}
\text{Gross investment:} & \quad I_t, \\
\text{Net investment:} & \quad I_t^{\text{net}} = I_t - \delta k_t = k_{t+1} - k_t, \\
\text{Depreciation:} & \quad D_t = \delta k_t.
\end{align}

\paragraph{The CES Constraint in Terms of Net Investment.}
Substituting $I = I^{\text{net}} + \delta k$ into the transformation:
\begin{equation}
\Hfun(C, I^{\text{net}} + \delta k) = A k.
\label{eq:ac-dep-ces-net}
\end{equation}

This shows that adjustment costs apply to \emph{gross} investment, not net investment. Even maintaining the capital stock ($I^{\text{net}} = 0$) requires resources $I = \delta k$, and these resources are subject to the transformation friction.

\paragraph{Interpretation.}
The CES transformation can be viewed as capturing:
\begin{itemize}[nosep]
\item \textbf{Installation costs}: Costs of putting new capital in place.
\item \textbf{Maintenance costs}: Costs of replacing depreciated capital.
\item \textbf{Sectoral reallocation}: Costs of shifting resources between consumption and investment goods sectors.
\end{itemize}

All three apply to gross investment flows, not just to changes in the capital stock.

\begin{calloutnote}[Why Gross Investment?]
The assumption that adjustment costs apply to gross investment (rather than net) is standard in the literature. The economic rationale is that the frictions---installation, learning, organizational change---occur whenever new capital goods are produced and installed, regardless of whether they replace depreciated capital or expand the stock.

An alternative specification would have adjustment costs on net investment: $\Hfun(C, k' - (1-\delta)k) = Y$. This would imply that replacement investment is frictionless, which may be appropriate in some contexts (e.g., identical replacement of worn-out machines) but not others (e.g., technological upgrading).
\end{calloutnote}

% --------------------------------------------------------------
\subsection{Comparison: Full vs. Incomplete Depreciation}
\label{subsec:ac-dep-comparison}
% --------------------------------------------------------------

\begin{table}[H]
\centering
\caption{Adjustment Costs with Full vs. Incomplete Depreciation}
\label{tab:ac-dep-comparison}
\renewcommand{\arraystretch}{1.6}
\begin{tabular}{@{}lcc@{}}
\toprule
\textbf{Object} & \textbf{Full ($\delta = 1$)} & \textbf{Incomplete ($\delta < 1$)} \\
\midrule
Capital accumulation & $k' = I$ & $k' = (1-\delta)k + I$ \\[4pt]
Return on capital & $R = \dfrac{A \, Q_t}{Q_{t+1}}$ & $R = \dfrac{A + (1-\delta) Q_{t+1}}{Q_t}$ \\[10pt]
Steady-state $Q$ & $\beta A$ & $\dfrac{\beta A}{1 - \beta(1-\delta)} = \dfrac{A}{r + \delta}$ \\[10pt]
Steady-state $I/Y$ & Determined by $C/I$ ratio & $\delta / A$ \\[4pt]
Net investment (SS) & $I$ & $0$ \\[4pt]
\bottomrule
\end{tabular}
\end{table}

% --------------------------------------------------------------
\subsection{Summary}
\label{subsec:ac-dep-summary}
% --------------------------------------------------------------

\begin{table}[H]
\centering
\caption{Incomplete Depreciation with CES Adjustment Costs}
\label{tab:ac-dep-summary}
\renewcommand{\arraystretch}{1.6}
\begin{tabular}{@{}ll@{}}
\toprule
\textbf{Object} & \textbf{Formula} \\
\midrule
Capital accumulation & $k' = (1-\delta) k + I$ \\[6pt]
Return on capital & $R_{t+1} = \dfrac{A + (1-\delta) Q_{t+1}}{Q_t}$ \\[10pt]
Euler equation & $u'(C_t) = \beta R_{t+1} \, u'(C_{t+1})$ \\[6pt]
Steady-state $Q$ & $Q^{ss} = \dfrac{A}{r + \delta}$ where $r = 1/\beta - 1$ \\[10pt]
User cost of capital & $(r + \delta) Q = A$ \\[6pt]
Steady-state investment & $I^{ss} = \delta k^{ss}$ (replacement only) \\[6pt]
\bottomrule
\end{tabular}
\end{table}

\begin{callouttip}[Key Insights]
Incomplete depreciation introduces:
\begin{itemize}[nosep]
\item \textbf{Capital gains in returns}: The return includes the resale value $(1-\delta) Q_{t+1}$ of undepreciated capital.
\item \textbf{User cost of capital}: Steady-state $Q = A/(r + \delta)$ reflects both interest and depreciation costs.
\item \textbf{Gross vs. net investment}: Adjustment costs apply to gross investment; even replacement investment is costly.
\item \textbf{Richer dynamics}: The state variable $k_t$ evolves gradually, creating transitional dynamics even with constant productivity.
\item \textbf{Steady-state investment}: In steady state, gross investment equals depreciation ($I = \delta k$) and net investment is zero.
\end{itemize}
The CES transformation continues to govern the consumption-investment trade-off, with Tobin's $Q$ summarizing the shadow cost of capital accumulation.
\end{callouttip}


